\begin{comment}
	Il programma si compone di diversi argomenti teorici, a continuazione del modulo di sistemi operativi in Elementi di Architettura e Sistemi Operativi, tutti corredati da corrispondenti esercizi. Tuttavia una parte considerevole del corso è orientata al laboratorio, al fine di studiare librerie POSIX per la sincronizzazione e sperimentare sul sistema operativo didattico MentOS alcuni dei concetti imparati a lezione.
	Gli argomenti di teoria sono i seguenti:
	- Schedulazione dei processi
	- Sincronizzazione fra processi e stallo
	- Sistema memoria: paginazione e caching
	- Periferiche di ingresso ed uscita
	- Sistema di file
	A questi argomenti corrispondono esercizi svolti in classe.
	Gli argomenti di laboratorio invece sono:
	- Meccanismi base di sincronizzazione e comunicazione inter/intra-processo: pipe e FIFO
	- Meccanismi di sincronizzazione e comunicazione inter-processo: l'IPC System V (semafori, memoria condivisa, coda di messaggi)
	- La libreria Pthread per i monitor dei processi
	- Il sistema operativo didattico MentOS: fondamentali, gestione dei processi, gestione della memoria, chiamate di sistema
\end{comment}